\section{Gradient}


\begin{bookfigure}{Gradient vectors are perpendicular to level curves and point
toward the steepest local increase of the scalar field.}
\begin{tikzpicture}[scale=1.0]
  \foreach \r/\lab in {0.8/1,1.5/2,2.25/3,3.05/4}{
    \draw[visualcyan!75,line width=0.9pt] (0,0) ellipse ({1.35*\r} and {\r});
    \node[visual label] at ({1.35*\r},0.2) {$f=\lab$};
  }
  \foreach \a in {20,65,115,165,215,265,315}{
    \draw[field vector] ({1.3*cos(\a)},{0.95*sin(\a)})
      -- ({2.25*cos(\a)},{1.65*sin(\a)});
  }
  \fill[visualgold] (1.15,0.52) circle (2pt);
  \draw[force vector] (1.15,0.52) -- (2.4,1.08)
    node[above right,visual label] {$\nabla f$};
  \node[font=\small,align=center,text width=4.2cm] at (7.2,0.4)
    {Level curves connect points of equal field value.\\[2mm]
     The gradient crosses them orthogonally.};
\end{tikzpicture}
\end{bookfigure}


If

\[
f(x,y,z)
\]

is a scalar field,

then

\[
\boxed{
	\nabla f
	=
	\left(
	\frac{\partial f}{\partial x},
	\frac{\partial f}{\partial y},
	\frac{\partial f}{\partial z}
	\right).
}
\]

The gradient points toward the direction of greatest increase of the function.

For example,

\[
T(x,y)
\]

may represent temperature.

Then

\[
\nabla T
\]

points in the direction where the temperature rises most rapidly.

%%%%%%%%%%%%%%%%%%%%%%%%%%%%%%%%%%%%%%%%%%%%%%%%%%%%%%%%%%%%%%%%%%%%%%%%%%%%%%%%


\section{The Gradient --- Measuring the Direction of Maximum Change}
This section introduces the gradient as the first differential operator. It maps a scalar field to a vector field indicating the direction of maximum increase.

\subsection{Motivation}
Consider a mountain landscape. The gradient points in the direction of steepest ascent and its magnitude gives the steepness.

\subsection{Definition}
For a scalar field $f(x,y,z)$,
\[
\nabla f=\frac{\partial f}{\partial x}\mathbf e_x+\frac{\partial f}{\partial y}\mathbf e_y+\frac{\partial f}{\partial z}\mathbf e_z.
\]
The operator $\nabla$ is called the nabla operator.

\subsection{Geometric interpretation}
The gradient:
\begin{itemize}
\item points toward the greatest increase,
\item has magnitude equal to the maximum rate of change,
\item is perpendicular to level surfaces.
\end{itemize}

\subsection{Examples}
For $f=x^2+y^2$, $\nabla f=2x\mathbf e_x+2y\mathbf e_y$.
Heat conduction obeys $q=-k\nabla T$.
Electrostatics obeys $\mathbf E=-\nabla V$.

\subsection{Optimization}
Gradient descent updates parameters by
\[
x_{\rm new}=x-\eta\nabla f.
\]

\subsection{Looking Ahead}
Gradient, divergence and curl form the core of vector calculus and underpin Maxwell's equations, the Schr\"odinger equation, fluid mechanics and field theory.
